\section{A target-mass bound at every width}\label{sec:target-mass}
Throughout this section $\nu$ has the even density in
\eqref{eq:curvature}. Adding a constant to $V$ changes nothing, so take
$V(0)=0$. Evenness gives $V'(0)=0$ and hence
$m t^2/2\leq V(t)\leq Mt^2/2$.

\begin{lemma}[Concentration and translation]\label{lem:concentration}
A product of $\nu$ coordinates, conditioned on any centrally symmetric
convex set of positive probability, is centered and satisfies
\begin{equation}\label{eq:conditional-subgaussian}
 \E e^{\lambda u\cdot Z}\leq
        \exp\{\lambda^2\|u\|_2^2/(2m)\}.
\end{equation}
For the unconditioned product measure $\mu$, Euclidean enlargement obeys
\begin{equation}\label{eq:enlargement}
 \mu(A_r)\geq\Phi\bigl(\Phi^{-1}(\mu(A))+\sqrt m\,r\bigr).
\end{equation}
Finally, if $E$ is any centrally symmetric event of positive probability
and $a$ is deterministic, then
\begin{equation}\label{eq:translation-mass}
 \mu(E+a)\geq\mu(E)\exp\{-M\|a\|_2^2/2\}.
\end{equation}
\end{lemma}
\begin{proof}
Conditioning on a convex set preserves $m$-strong log-concavity. The
Brascamp--Lieb variance inequality \cite{brascamp-lieb}, applied after
any linear exponential tilt, bounds the Hessian of its log moment
generating function by $m^{-1}I$. Central symmetry makes the first
derivative at zero vanish. Integrating twice proves
\eqref{eq:conditional-subgaussian}. Convex indicator potentials follow
by approximation with smooth convex barriers.

For \eqref{eq:enlargement}, the Caffarelli contraction theorem
\cite{caffarelli} gives an $m^{-1/2}$-Lipschitz transport from standard
Gaussian measure to $\mu$. The inverse image of $A_r$ contains the
$\sqrt m\,r$ enlargement of the inverse image of $A$. Gaussian
isoperimetry \cite{ledoux} gives the displayed inequality.

For the last assertion let $Q$ be the law of $Z+a$ conditional on
$Z\in E$. Direct calculation of relative entropy gives
\[
 D(Q\Vert\mu)=-\log\mu(E)
 +\E\left[\sum_i(V(Z_i+a_i)-V(Z_i))\mid E\right].
\]
Taylor's theorem bounds the sum by
$\sum_i a_i V'(Z_i)+M\|a\|_2^2/2$. The conditional expectation of
its linear term is zero because $V'$ is odd and the conditional law is
invariant under simultaneous negation. Since $Q$ is supported on $E+a$,
$D(Q\Vert\mu)\geq-\log\mu(E+a)$. This proves
\eqref{eq:translation-mass}.
\end{proof}

\begin{lemma}[One-dimensional small balls]\label{lem:small-ball}
Every centered even log-concave variable $Y$ of variance at most $1/m$
satisfies
\[
 \Pr\{|Y|\leq\rho\}\geq c_m\rho\quad(0<\rho\leq1),
 \qquad c_m=\tfrac38\min\{1,\sqrt m\}.
\]
This applies to every normalized sum of disjoint coordinates of $Z$.
\end{lemma}
\begin{proof}
Let $R=2\max\{1,m^{-1/2}\}$. Chebyshev gives probability at least
$3/4$ in $[-R,R]$. Symmetric unimodality gives probability at least
$(\rho/R)(3/4)$ in $[-\rho,\rho]$. Log-concavity is preserved by
convolution and linear projection, and normalized independent sums retain
variance at most $1/m$.
\end{proof}

We first construct an exact clone of a target. Partition the $N$ neurons
of each hidden layer into $C$ groups of sizes $k_a\in\{\lfloor N/C\rfloor,
\lceil N/C\rceil\}$. Copy the corresponding teacher neuron throughout
each group. For an internal teacher weight $W_{ab}$ use $W_{ab}/k_b$
on every edge from group $b$ to group $a$. Divide the readout weight
from group $a$ by $k_a$. The cloned network computes exactly the same
function. Let $P$ project onto the group-constant subspace, with orthonormal
basis $u_a=k_a^{-1/2}\mathbf1_{\mathrm{group}\ a}$.

Define the weighted energy
\begin{equation}\label{eq:kappa}
 \kappa_N(w)=\frac dN\|W^1\|_F^2+
     \sum_{s=2}^{L-1}\|W^s\|_F^2+N\|W^L\|_2^2,
 \qquad \kappa_N(g)=\inf_{w\in U_{N,g}}\kappa_N(w).
\end{equation}
The first-layer clone contribution is at most $d(d+1)C^2$.
Because $\max k_a/\min k_a\leq2$ and $N/\min k_a\leq2C$,
its internal contribution is at most $2(L-2)C^4$ and its readout
contribution at most $2C^4$. Amplify the teacher readout by three; the
resulting exact clone has signed margin greater than $3/2$. A convenient
common bound for its energy is
\begin{equation}\label{eq:B0}
 B_0=d(d+1)C^2+2(L-2)C^4+36C^4.
\end{equation}
The larger final constant leaves harmless slack.

\begin{proposition}[Target mass]\label{prop:target-mass}
For each fixed $L$ there is a polynomially bounded function $A_L(d,C)$,
depending also on $m,M$, such that
\begin{equation}\label{eq:target-mass}
 \mu_{N,\gamma}(U_{N,f})\geq e^{-A_L(d,C)N^\gamma}
 \qquad(\gamma\geq1,\ N\geq C=C_L(f)).
\end{equation}
\end{proposition}
\begin{proof}
Put $t=N^{(1-\gamma)/2}\leq1$. In standardized coordinates impose the
centered cylinder $\mathcal C_\rho$ consisting of the following constraints:
all first-layer and readout coordinates have absolute value at most $\rho$;
each internal row obeys $|\langle Z_{j,*},u_a\rangle|\leq\rho$ for
every group $a$. These constraints involve independent random variables:
within an internal row the sums concern disjoint groups. There are
\[
 r_0=N[d+2+(L-2)C]
\]
constraints. Lemma~\ref{lem:small-ball} gives
$\mu(\mathcal C_\rho)\geq(c_m\rho)^{r_0}$.

Conditional on the cylinder, every internal matrix is centered and
$m$-strongly log-concave. For unit vectors $u,v$, the scalar $u^TZv$
is therefore $m^{-1/2}$-sub-Gaussian. Two $1/4$-nets, each with at most
$9^N$ elements, give
\[
 \Pr\{\|Z\|_{\mathrm{op}}>2q\sqrt N\mid\mathcal C_\rho\}
 \leq2\exp\{(2\log9-mq^2/2)N\}.
\]
Take
\[
 b_{m,L}=2\sqrt{\frac2m(2\log9+\log(4L))}.
\]
Each conditional failure probability for
$\|Z\|_{\mathrm{op}}\leq b_{m,L}\sqrt N$ is at most $1/(2L)$.
Let $E$ be the cylinder intersected with these bounds for all internal
layers. It is centrally symmetric and has
$\mu(E)\geq\tfrac12(c_m\rho)^{r_0}$.

Translate this event by the amplified clone in standardized coordinates.
The internal weight perturbation decomposes as
\[
 \frac t{\sqrt N}Z=R+E_\perp,
 \qquad R=\frac t{\sqrt N}ZP,
 \qquad E_\perp=\frac t{\sqrt N}Z(I-P),
\]
where
$\|R\|_{\mathrm{op}}\leq\rho t\sqrt C$,
$E_\perp P=0$, and $\|E_\perp\|_{\mathrm{op}}\leq b_{m,L}t$.
We detail stability to make the role of the cylinder explicit. Write
$e_s=N^{-1/2}\|h^s-h_*^s\|_2$, with $h_*^s$ the exact clone's
activations. Its internal matrices have norm at most $\sqrt2C^2$, and
$h_*^s\in\mathrm{range}(P)$. The first layer has
\[
 e_1\leq \rho t(d+1)/\sqrt d\leq\rho(d+1).
\]
For $s\geq2$, boundedness and the contraction property of tanh give
\[
 e_s\leq(\sqrt2C^2+b_{m,L})e_{s-1}+\rho\sqrt C.
\]
Set $A_0=2+3C^2+b_{m,L}$. Induction bounds
$e_{L-1}\leq\rho(d+L+1)A_0^{L-2}$.
The cloned readout has $\sqrt N\|v_*\|_2\leq3\sqrt2 C^2$,
and the perturbation readout has $\ell^1$ norm at most $\rho t\leq\rho$.
Consequently the output error is at most
\[
 3\sqrt2 C^2\rho(d+L+1)A_0^{L-2}+\rho.
\]
Choose
\[
 \rho=[2+9C^2(d+L+1)A_0^{L-2}]^{-1}.
\]
The last error is strictly smaller than $1/2$, so the translated event
lies in $U_{N,f}$, indeed with signed margin greater than one.

If $a$ is the standardized translation, then
$\|a\|_2^2=N^\gamma\kappa_N(w_*)\leq B_0N^\gamma$.
Lemma~\ref{lem:concentration} gives
\[
 -\log\mu_{N,\gamma}(U_{N,f})
 \leq\tfrac M2B_0N^\gamma
 +N[d+2+(L-2)C]\log\frac1{c_m\rho}+\log2.
\]
Since $N\leq N^\gamma$, one may take
\begin{equation}\label{eq:Aexplicit}
 A_L(d,C)=1+\tfrac M2B_0+
 [d+2+(L-2)C]\log\frac1{c_m\rho}+\log2.
\end{equation}
At fixed depth and fixed curvature constants this has a polynomial
majorant in $d,C$.
\end{proof}
